\section{Conclusão}

A manutenção do engajamento em atividades físicas representa um desafio contínuo, frequentemente prejudicado pela monotonia e pela ausência de um retorno interativo imediato \cite{plangger2022little}. Este trabalho apresentou uma proposta de arquitetura estruturada entre \textit{smartwatch}, aplicativo móvel e \textit{backend} para fornecer uma experiência de gamificação imersiva. A avaliação de desempenho preliminar sugere que os tempos de resposta ponta a ponta observados ficaram inferiores a 120 ms, compatíveis com a sincronização da mecânica de fuga em ciclos de aproximadamente um segundo.

A relevância desta solução reside na superação das limitações das abordagens comerciais genéricas. Ao invés de oferecer um acompanhamento puramente reativo, a integração de métricas biométricas em tempo real possibilita um ambiente responsivo focado em permanecer o usuário dentro de seu canal de fluxo \cite{csikszentmihalyi1989optimal}, visando fortalecer a motivação intrínseca de forma personalizada \cite{cernev2012autodeterminacao}.

Como trabalhos futuros e extensões da plataforma, sugere-se a implementação da funcionalidade de pausa durante a sessão, permitindo maior flexibilidade em cenários reais de treino. Além disso, a incorporação da leitura de ganho de altitude enriqueceria a coleta de telemetria, impactando diretamente o cálculo do esforço físico na corrida. A adição de efeitos sonoros no fone de ouvido pode aumentar a imersão, reforçando a sensação de perseguição pela horda e tornando a experiência mais realista.

A adaptação automática de dificuldade também é indicada como uma evolução futura do sistema. Embora tenha sido prevista conceitualmente, sua implementação completa não foi incorporada nesta etapa devido à dificuldade de garantir precisão suficiente na variação dos dados de telemetria durante as validações técnicas. Para obter uma coleta de resultados mais frequente, estável e coerente, optou-se por manter a lógica da horda em um modelo mais controlado, deixando os ajustes dinâmicos de dificuldade para uma etapa posterior.

No âmbito técnico, outras melhorias incluem o desenvolvimento de paridade do fluxo de telemetria e renderização para a plataforma iOS, a aplicação de mecanismos de suavização e limites de ritmo nas hordas adaptativas, bem como a validação do sistema com sensores físicos em cenários automatizados. Também se propõe, como continuidade, a implementação de rankeamento das infestações entre contas conectadas, integração com listas de amigos adicionados e aprimoramento da visualização da gameplay. Atualmente, a cena visual atua principalmente como uma demonstração do processo de corrida e perseguição. Futuramente, o zumbi poderá reagir visualmente à distância em relação ao personagem, de forma semelhante ao que já pode ser acompanhado pela barra de progresso exibida tanto no \textit{smartwatch} quanto no aplicativo móvel.